\section{Evaluation Methodology}

% Scope: introduce the bottom-up framework and point to the benchmark
% companion paper. Keep brief; results sections follow.

\subsection{Bottom-Up Component Evaluation}
% Why retriever first, then generator under oracle context, then end-to-end.
% Coupling problem: a weak retriever makes generator eval meaningless;
% end-to-end alone confounds retrieval quality with generation faithfulness.

\subsection{Benchmarks Used}
% mamaretrieval (~3,185 queries, ~78,571 labeled query-chunk pairs):
% retriever evaluation and oracle context for generator isolation.
% mamabench (~25,997 rows across 7 sources, 3 set_types): end-to-end
% MCQ + open-ended evaluation. Construction details in the companion
% benchmarks paper [cite paper2].

\subsection{Metric Choices}
% Retrieval: Hit Rate, MRR, nDCG, Recall, Precision. Lead with
% completeness-robust metrics; report Recall/Precision with an
% audit-gap error bar [cite paper2 \S Completeness Audit].
% Generation under oracle: MiniCheck sentence-level support rate
% (Pipeline 1 + Pipeline 2), LLM-judge calibration, stability probes.
% End-to-end: MCQ accuracy + calibration (Brier, ECE); open-ended
% rubric scores (factuality / reasoning / harm / omission / guideline
% adherence); safety pass rates and robustness.
% Latency: P50 / P90 / P99 percentiles + full distribution plots.
